\documentclass[11pt,a4paper]{article}

% =========================================================
% COMMON PACKAGES
% =========================================================
\usepackage{lmodern}
\usepackage[T1]{fontenc}
\usepackage[utf8]{inputenc}
\usepackage[english]{babel}

\usepackage{amsmath,amssymb,amsfonts,amsthm,mathtools}
\usepackage{bm}
\usepackage{microtype}
\usepackage{enumitem}
\usepackage{csquotes}
\usepackage{mathrsfs}
\usepackage{thmtools}
\usepackage{hyperref}
\usepackage[nameinlink,capitalize]{cleveref}
\usepackage{geometry}

\geometry{
	a4paper,
	margin=2.8cm
}

% =========================================================
% THEOREM STYLES
% =========================================================
\numberwithin{equation}{section}

\declaretheoremstyle[
headfont=\bfseries,
bodyfont=\itshape,
spaceabove=8pt,
spacebelow=8pt
]{plainstyle}

\declaretheoremstyle[
headfont=\bfseries,
bodyfont=\normalfont,
spaceabove=8pt,
spacebelow=8pt
]{defstyle}

\declaretheoremstyle[
headfont=\itshape,
bodyfont=\normalfont,
spaceabove=4pt,
spacebelow=4pt
]{remarkstyle}

\declaretheorem[style=plainstyle,name=Theorem]{theorem}
\declaretheorem[style=plainstyle,name=Proposition,sibling=theorem]{proposition}
\declaretheorem[style=plainstyle,name=Lemma,sibling=theorem]{lemma}
\declaretheorem[style=plainstyle,name=Corollary,sibling=theorem]{corollary}

\declaretheorem[style=defstyle,name=Definition,sibling=theorem]{definition}
\declaretheorem[style=defstyle,name=Example,sibling=theorem]{example}

\declaretheorem[style=remarkstyle,name=Remark,sibling=theorem]{remark}

% Cleveref names
\crefname{theorem}{Theorem}{Theorems}
\crefname{proposition}{Proposition}{Propositions}
\crefname{lemma}{Lemma}{Lemmas}
\crefname{corollary}{Corollary}{Corollaries}
\crefname{definition}{Definition}{Definitions}
\crefname{example}{Example}{Examples}
\crefname{remark}{Remark}{Remarks}
\crefname{equation}{Equation}{Equations}
\crefname{section}{Section}{Sections}

% =========================================================
% COMMON MACROS FOR THE TWIN-SPACES SERIES
% =========================================================

% Sets and fields
\newcommand{\R}{\mathbb{R}}
\newcommand{\C}{\mathbb{C}}
\newcommand{\Z}{\mathbb{Z}}
\newcommand{\N}{\mathbb{N}}

% Operators
\DeclareMathOperator{\End}{End}
\DeclareMathOperator{\Hom}{Hom}
\DeclareMathOperator{\id}{id}
\DeclareMathOperator{\Ad}{Ad}
\DeclareMathOperator{\Aut}{Aut}
\DeclareMathOperator{\Ker}{Ker}
\DeclareMathOperator{\Img}{Im}
\DeclareMathOperator{\HH}{HH}

% Twin-space notation
\newcommand{\TwinSpace}{(E,\alpha,G)}
\newcommand{\Twin}{\mathcal{T}}
\newcommand{\TwinOrbit}{O}
\newcommand{\LinearSubgroup}{L}

% Cohomology notation
\newcommand{\Cochain}{C}
\newcommand{\diff}{\delta}

% Misc
\newcommand{\abs}[1]{\left\lvert #1 \right\rvert}
\newcommand{\norm}[1]{\left\lVert #1 \right\rVert}
\newcommand{\bracket}[1]{\left\langle #1 \right\rangle}

% =========================================================
% METADATA FOR ARTICLE B
% =========================================================

\title{\textbf{Twin Spaces as 2-Categorical and Operadic Objects}}
\author{Jocelyn Nembe}
\date{\today}

% =========================================================
% DOCUMENT
% =========================================================

\begin{document}
	
	\maketitle
	
	\begin{abstract}
		In Article~I we introduced \emph{twin operations} attached to a group action, and in Article~A we developed 
		a deformation cohomology theory for twin spaces via a Hochschild-type complex $\Cochain^\bullet(G,\End(E))$.
		The goal of the present article is to show that twin spaces carry a rich higher-structural content:
		they naturally form a 2-category, admit a monoidal structure, and can be organized into operads whose algebras
		encode classical and new symmetry-compatible algebraic structures.
		
		More precisely, we define a 2-category $\Twin$ whose objects are twin spaces $\TwinSpace$, 
		whose 1-morphisms are equivariant structure-preserving maps, and whose 2-morphisms are natural
		transformations compatible with the twin operations.
		We construct a tensor product making $\Twin$ into a monoidal 2-category.
		We then introduce a family of operads $\mathcal{O}_G$ of twin operations associated with a group $G$,
		and show that many standard structures (associative $G$-algebras, $G$-equivariant module categories, 
		and certain crossed product constructions) appear as algebras over these operads.
		
		Conceptually, this article clarifies the categorical level at which twin spaces live:
		they are not just sets with operations, but 2-categorical and operadic objects whose deformation theory 
		(Article~A) can be viewed as the first step towards a higher, possibly derived, theory of twin structures.
	\end{abstract}
	
	\tableofcontents
	
	%===========================================================
	\section{Introduction}
	\label{sec:intro}
	
	Twin operations were introduced in Article~I as a new way to attach algebraic structure to a group action:
	given a set $E$ with a left action of a group $G$ and a binary operation $\alpha : E\times E\to E$, one considers
	the family of \emph{twin operations}
	\begin{equation}\label{eq:twin-def-B}
		\alpha_g(x,y) := g^{-1}\cdot \alpha(g\cdot x,\ g\cdot y), \qquad g\in G.
	\end{equation}
	The set $\TwinOrbit=\{\alpha_g : g\in G\}$ of such operations is canonically identified with the homogeneous space
	$\LinearSubgroup\backslash G$ (right cosets; see \cref{prop:orbit}), where $\LinearSubgroup=\{g\in G : \alpha_g=\alpha\}$ is the linear subgroup.
	A triple $\TwinSpace$ with this orbit structure was called a \emph{twin space}.
	
	In Article~A, we showed that twin spaces admit a deformation cohomology theory, based on a Hochschild-type
	complex $\Cochain^\bullet(G,\End(E))$ whose cohomology $\HH^\bullet(G,\End(E))$ controls $G$-equivariant
	deformations and rigidity of the twin structure.
	
	The aim of the present article is to study twin spaces from the \emph{higher-structural} point of view.
	In particular:
	
	\begin{itemize}
		\item we construct a 2-category $\Twin$, with twin spaces as objects, equivariant morphisms as 1-morphisms,
		and compatible natural transformations as 2-morphisms;
		\item we equip $\Twin$ with a tensor product making it a monoidal 2-category in the sense of 
		higher category theory~\cite{MacLane1998,Leinster2004,Kelly1982};
		\item for a fixed group $G$, we organize $G$-equivariant $n$-ary twin operations into an operad
		$\mathcal{O}_G$, in the sense of May and Loday--Vallette~\cite{May1972,LodayVallette2012}, and show that various
		algebraic structures correspond to algebras over $\mathcal{O}_G$;
		\item we interpret some aspects of the deformation theory of twin spaces (Article~A) as living inside
		this 2-categorical and operadic environment.
	\end{itemize}
	
	Conceptually, the message is that twin spaces form not just a set-theoretic or algebraic construction,
	but fit naturally into the framework of higher category theory and operads.
	This clarifies the role of twin spaces as a \emph{foundational} object for subsequent work on 
	derived, homotopical and quantum generalizations.

	There are three independent reasons to expect this higher structure, and the article makes each precise.
	First, the very definition of a twin space produces not one operation but an entire orbit
	$\TwinOrbit\cong\LinearSubgroup\backslash G$, and an orbit of objects always asks for morphisms relating its
	points---so a merely $1$-categorical packaging is incomplete. Second, the deformation theory of Article~A is
	controlled by a differential graded Lie algebra (its Maurer--Cartan equation), and DGLAs are the infinitesimal
	shadow of operadic and higher-categorical structure. Third, twin operations of higher arity compose by
	substitution exactly as the operations of an operad do. The $2$-category, the monoidal structure, and the
	operad constructed below are the precise homes for these three observations.
	
	The article is organized as follows.  
	In \cref{sec:2cat}, we define the 2-category $\Twin$ of twin spaces.  
	In \cref{sec:monoidal}, we construct a monoidal structure on $\Twin$.  
	In \cref{sec:operads}, we introduce the operads of twin operations $\mathcal{O}_G$ and describe their algebras.
	In \cref{sec:functoriality}, we discuss functoriality and 2-functors relating $\Twin$ to more classical categories.
	In \cref{sec:interaction}, we sketch how the deformation cohomology of Article~A can be seen from a 2-categorical/operadic point of view.
	Finally, in \cref{sec:examples} we illustrate the constructions on concrete examples.
	
	%===========================================================
	\section{Twin Spaces as Objects of a 2-Category}
	\label{sec:2cat}
	
	We begin by recalling the basic definition of a twin space and then organize twin spaces into a 2-category.
	
	\subsection{Twin spaces}
	
	Let $G$ be a group acting on a set $E$ from the left.
	Given a binary operation $\alpha:E\times E\to E$, we define the family of twin operations
	by \eqref{eq:twin-def-B}. Write again
	\[
	\TwinOrbit = \{\alpha_g : g\in G\}.
	\]

	\begin{proposition}[Orbit identification]\label{prop:orbit}
		The map $g\mapsto\alpha_g$ has fibres the \emph{right} cosets of
		$\LinearSubgroup=\{g\in G:\alpha_g=\alpha\}$: indeed $\alpha_{g_1}=\alpha_{g_2}$
		iff $g_1 g_2^{-1}\in\LinearSubgroup$, i.e.\ iff $\LinearSubgroup g_1=\LinearSubgroup g_2$.
		Hence $\TwinOrbit\cong\LinearSubgroup\backslash G$ (right cosets); when
		$\LinearSubgroup\trianglelefteq G$ this coincides with $G/\LinearSubgroup$.
	\end{proposition}
	
	\begin{definition}[Twin space]
		A \emph{twin space} is a triple $\TwinSpace$, where:
		\begin{itemize}
			\item $E$ is a set (or object in a fixed ambient category, e.g.\ a vector space),
			\item $G$ is a group acting on $E$,
			\item $\alpha:E\times E\to E$ is a binary operation such that the orbit
			$\TwinOrbit$ of \eqref{eq:twin-def-B} is well-defined and carries the induced $G$-action.
		\end{itemize}
		The associated \emph{linear subgroup} is $\LinearSubgroup := \{g\in G : \alpha_g=\alpha\}$.
	\end{definition}
	
	In many settings, one requires $\alpha$ to satisfy extra identities (associativity, commutativity, etc.),
	but for the 2-categorical structure the basic notion suffices.
	
	\subsection{1-morphisms: equivariant structure-preserving maps}
	
	We want morphisms that preserve both the $G$-action and the twin operation structure.
	
	\begin{definition}[1-morphisms of twin spaces]
		Let $\TwinSpace$ and $(E',\alpha',G')$ be twin spaces.
		A \emph{1-morphism} 
		\[
		F : (E,\alpha,G)\longrightarrow (E',\alpha',G')
		\]
		consists of:
		\begin{itemize}
			\item a group homomorphism $\phi:G\to G'$,
			\item a map of sets (or linear map) $f:E\to E'$,
		\end{itemize}
		such that:
		\begin{enumerate}[label=(\alph*)]
			\item $f$ is $\phi$-equivariant:
			\[
			f(g\cdot x) = \phi(g)\cdot f(x), \qquad \forall g\in G,\ x\in E;
			\]
			\item $f$ is compatible with the operations:
			\[
			f(\alpha(x,y)) = \alpha'(f(x),f(y)), \qquad \forall x,y\in E.
			\]
		\end{enumerate}
	\end{definition}
	
	\begin{remark}
		Condition (b) can be strengthened if needed to require compatibility with \emph{all} twin operations
		$\alpha_g$ and $\alpha'_{\phi(g)}$, but this already follows from (a) and the definition of twin operations:
		\[
		f(\alpha_g(x,y))
		= f\big(g^{-1}\cdot \alpha(g\cdot x,g\cdot y)\big)
		= \phi(g)^{-1}\cdot f(\alpha(g\cdot x,g\cdot y))
		= \phi(g)^{-1}\cdot \alpha'(f(g\cdot x),f(g\cdot y)),
		\]
		and $f(g\cdot x)=\phi(g)\cdot f(x)$ implies
		\[
		f(\alpha_g(x,y)) = (\alpha')_{\phi(g)}(f(x),f(y)).
		\]
	\end{remark}
	
	\begin{example}[The determinant as a $1$-morphism]\label{ex:det-morphism}
		Let $\TwinSpace$ be the matrix twin space $E=M_n(k)$, $\alpha(X,Y)=XY$, with $G=GL_n(k)$ acting by
		conjugation, and let $(E',\alpha',G')=(k,\,\cdot\,,\{1\})$ be the multiplicative twin space on the ground
		field with trivial group. The pair $F=(\det,\phi)$, with $f=\det:M_n(k)\to k$ and $\phi:GL_n(k)\to\{1\}$
		the trivial homomorphism, is a $1$-morphism: equivariance holds because
		$\det(gXg^{-1})=\det X=\phi(g)\cdot\det X$, and operation-compatibility is the multiplicativity
		$\det(XY)=\det X\,\det Y$. Thus the determinant is not merely a function but a morphism of twin spaces,
		collapsing the conjugation-invariant matrix multiplication onto ordinary multiplication of scalars.
	\end{example}

	\subsection{2-morphisms: compatible natural transformations}
	
	We now define 2-morphisms between 1-morphisms.
	
	\begin{definition}[2-morphisms of twin spaces]
		Let $F=(f,\phi)$ and $F'=(f',\phi')$ be 1-morphisms from $\TwinSpace$ to $(E',\alpha',G')$.
		A \emph{2-morphism}
		\[
		\eta : F \Rightarrow F'
		\]
		is a family of maps
		\[
		\eta_x : f(x) \to f'(x), \qquad x\in E,
		\]
		in the ambient category (typically $\id_{f(x)}$ if one stays within sets or vector spaces),
		such that:
		\begin{enumerate}[label=(\alph*)]
			\item $\eta$ is $G$-equivariant:
			\[
			\eta_{g\cdot x} = \phi(g)\cdot \eta_x, \qquad \forall g\in G,\ x\in E,
			\]
			which is well-typed because $f(g\cdot x)=\phi(g)\cdot f(x)$ and $f'(g\cdot x)=\phi(g)\cdot f'(x)$, so that the
			translate $\phi(g)\cdot\eta_x$ is again a morphism $f(g\cdot x)\to f'(g\cdot x)$;
			\item $\eta$ is compatible with the operations:
			\[
			\eta_{\alpha(x,y)} = \alpha'\big(\eta_x,\eta_y\big),
			\]
		\end{enumerate}
	\end{definition}
	
	\begin{remark}
		In the purely set-theoretic case, the above notion collapses to the identity transformation
		(i.e.\ 2-morphisms are trivial). The main interest is in enriched contexts, e.g.\ when $E$ is a category
		or a module object, in which case $\eta_x$ is a genuine morphism in that category. In that enriched setting the
		two conditions read transparently: equivariance says the component at $g\cdot x$ is the $\phi(g)$-translate of
		the component at $x$, while operation-compatibility says that evaluating $\eta$ at $\alpha(x,y)$ agrees with
		applying the operation $\alpha'$, viewed as a bifunctor, to the pair $(\eta_x,\eta_y)$. When the ambient
		category is $\mathsf{Set}$ or $\mathsf{Vect}$ the only available components force $f=f'$ and $\eta=\id$, which
		is exactly why the genuinely $2$-categorical content surfaces only after enrichment (for instance in
		$\mathsf{Cat}$, where a $2$-morphism is a $G$-equivariant natural transformation between equivariant functors).
		For clarity, we adopt the 2-categorical language even though in some basic examples the 2-morphisms are only identities.
	\end{remark}
	
	\subsection{The 2-category \texorpdfstring{$\Twin$}{Twin}}
	
	\begin{theorem}\label{thm:Twin-2cat}
		Twin spaces, 1-morphisms and 2-morphisms as above form a (strict) 2-category $\Twin$.
	\end{theorem}
	
	\begin{proof}
		The composition of 1-morphisms is defined componentwise:
		if $F=(f,\phi):\TwinSpace\to (E',\alpha',G')$ and $F'=(f',\phi'):(E',\alpha',G')\to (E'',\alpha'',G'')$,
		then
		\[
		F'\circ F = (f'\circ f,\ \phi'\circ \phi).
		\]
		Equivariance and compatibility with operations are preserved by composition.
		The identity 1-morphism is $(\id_E,\id_G)$.
		
		The vertical and horizontal compositions of 2-morphisms are defined in the usual way for natural
		transformations, and the coherence axioms of a 2-category are satisfied by standard arguments
		(see e.g.~\cite{MacLane1998,Leinster2004}). The middle-four interchange law holds because it holds in the
		ambient enriched category, and both the equivariance condition $\eta_{g\cdot x}=\phi(g)\cdot\eta_x$ and the
		operation-compatibility condition are visibly closed under vertical and horizontal composition. Thus we obtain a strict 2-category.
	\end{proof}
	
	%===========================================================
	\section{Monoidal Structure on the 2-Category \texorpdfstring{$\Twin$}{Twin}}
	\label{sec:monoidal}
	
	We now equip $\Twin$ with a tensor product, making it into a monoidal 2-category.
	
	\subsection{Tensor product of twin spaces}
	
	Let $\TwinSpace$ and $(E',\alpha',G')$ be twin spaces.
	
	\begin{definition}[Tensor product of twin spaces]
		The \emph{tensor product} of $\TwinSpace$ and $(E',\alpha',G')$ is the twin space
		\[
		(E\otimes E',\ \alpha\otimes\alpha',\ G\times G'),
		\]
		defined as follows:
		\begin{itemize}
			\item the underlying object is $E\otimes E'$ (e.g.\ the cartesian product $E\times E'$ for sets,
			the tensor product for vector spaces);
			\item the group $G\times G'$ acts componentwise:
			\[
			(g,g')\cdot (x,x') := (g\cdot x,\ g'\cdot x');
			\]
			\item the operation $\alpha\otimes\alpha'$ is given by
			\[
			(\alpha\otimes\alpha')\big((x_1,x'_1),(x_2,x'_2)\big)
			:=\big(\alpha(x_1,x_2),\ \alpha'(x'_1,x'_2)\big).
			\]
		\end{itemize}
	\end{definition}
	
	\begin{proposition}
		The triple $(E\otimes E',\alpha\otimes\alpha',G\times G')$ is a twin space, and twin operations satisfy
		\[
		(\alpha\otimes\alpha')_{(g,g')} = \alpha_g\otimes\alpha'_{g'}.
		\]
	\end{proposition}
	
	\begin{proof}
		One checks directly that
		\[
		(\alpha\otimes\alpha')_{(g,g')}\big((x_1,x'_1),(x_2,x'_2)\big)
		= (g,g')^{-1}\cdot (\alpha\otimes\alpha')\big((g,g')\cdot(x_1,x'_1),(g,g')\cdot(x_2,x'_2)\big)
		\]
		unwinds, using the componentwise action and the definition of $\alpha\otimes\alpha'$, as
		\[
			(g^{-1},g'^{-1})\cdot\big(\alpha(g\cdot x_1,g\cdot x_2),\,\alpha'(g'\cdot x'_1,g'\cdot x'_2)\big)
			=\big(g^{-1}\cdot\alpha(g\cdot x_1,g\cdot x_2),\ g'^{-1}\cdot\alpha'(g'\cdot x'_1,g'\cdot x'_2)\big)
			=\big(\alpha_g(x_1,x_2),\,\alpha'_{g'}(x'_1,x'_2)\big),
		\]
		which is precisely $(\alpha_g\otimes\alpha'_{g'})\big((x_1,x'_1),(x_2,x'_2)\big)$. That the triple is a twin
		space is then immediate: the orbit of $\alpha\otimes\alpha'$ is the product of the two orbits.
	\end{proof}
	
	\begin{remark}[Linear subgroup of a tensor product]\label{rem:tensor-L}
		The linear subgroup of the tensor product is the product of the linear subgroups,
		$\LinearSubgroup_{E\otimes E'}=\LinearSubgroup\times\LinearSubgroup'$, because
		$(\alpha\otimes\alpha')_{(g,g')}=\alpha\otimes\alpha'$ holds iff $\alpha_g=\alpha$ and $\alpha'_{g'}=\alpha'$.
		Hence the twin orbit of a tensor product factors as a product of orbits,
		\[
			(\LinearSubgroup\times\LinearSubgroup')\backslash(G\times G')\;\cong\;(\LinearSubgroup\backslash G)\times(\LinearSubgroup'\backslash G'),
		\]
		so the twin-orbit $2$-functor of \cref{sec:functoriality} is itself monoidal.
	\end{remark}

	\subsection{Tensor product of morphisms and monoidal axioms}
	
	\begin{definition}
		Given 1-morphisms $F=(f,\phi):\TwinSpace\to (E_1,\alpha_1,G_1)$
		and $F'=(f',\phi'):(E',\alpha',G')\to (E'_1,\alpha'_1,G'_1)$, we define
		\[
		F\otimes F' := (f\otimes f',\ \phi\times \phi'),
		\]
		where
		\[
		(f\otimes f')(x,x') := (f(x),f'(x')),\qquad
		(\phi\times\phi')(g,g') := (\phi(g),\phi'(g')).
		\]
		Tensor product of 2-morphisms is defined pointwise.
	\end{definition}
	
	\begin{theorem}\label{thm:Twin-monoidal}
		Equipped with the tensor product above, the 2-category $\Twin$ becomes a monoidal 2-category.
	\end{theorem}
	
	\begin{proof}
		Associativity and unit constraints are inherited from the ambient category (e.g.\ Sets or Vect),
		and the coherence axioms of a monoidal 2-category follow from standard arguments as in 
		\cite{MacLane1998,Leinster2004}.
	\end{proof}
	
	\begin{remark}
		In many concrete settings (e.g.\ vector spaces over a fixed field), the monoidal structure is in fact strict
		or can be strictified without loss of generality, so that associativity and unit isomorphisms can be taken
		to be identities.
	\end{remark}
	
	%===========================================================
	\section{Operads of Twin Operations}
	\label{sec:operads}
	
	We now introduce a family of operads whose operations are twin operations compatible with a given group action.
	
	\subsection{The underlying collections}
	
	Fix a group $G$ and a $G$-set (or $G$-object) $E$.
	
	\begin{definition}[Twin $n$-ary operations]
		For $n\ge 1$, let
		\[
		\mathcal{O}_G(n) := \{ \alpha : E^n\to E \mid \alpha \text{ is $G$-equivariant in the sense of twin operations} \}.
		\]
		More precisely, $\alpha\in\mathcal{O}_G(n)$ if there exists a family
		\[
		\alpha_g : E^n\to E,\qquad g\in G,
		\]
		such that
		\[
		\alpha_g(x_1,\dots,x_n) = g^{-1}\cdot \alpha(g\cdot x_1,\dots,g\cdot x_n),
		\]
		and the map $g\mapsto \alpha_g$ is compatible with the group action (in the same spirit as in the binary case).
	\end{definition}
	
	We think of $\mathcal{O}_G(n)$ as the space of $n$-ary twin operations on $E$.

	\begin{remark}[$\mathcal{O}_G$ is the endomorphism operad internal to $G$-sets]\label{rem:OG-internal}
		The existence of a transport family $g\mapsto\alpha_g$ is automatic: \emph{every} map $\alpha:E^n\to E$
		admits one, namely $\alpha_g(\mathbf{x})=g^{-1}\cdot\alpha(g\cdot\mathbf{x})$. Thus, set-theoretically,
		$\mathcal{O}_G(n)=\mathrm{Maps}(E^n,E)$ is the full endomorphism operad $\End_E(n)$; the twin viewpoint adds
		not more operations but more \emph{structure}, namely the $G$-action $g\cdot\alpha:=\alpha_{g^{-1}}$ on each
		$\End_E(n)$, compatible with substitution. In one phrase, $\mathcal{O}_G$ is the endomorphism operad
		\emph{internal to the category of $G$-sets}, and the linear subgroup $\LinearSubgroup$ of a chosen $\alpha$
		is simply its stabiliser under this action.
	\end{remark}
	
	\subsection{Operadic composition}
	
	Let $\alpha\in\mathcal{O}_G(k)$ and $\beta_1\in\mathcal{O}_G(n_1),\dots,\beta_k\in\mathcal{O}_G(n_k)$.
	We define their \emph{substitution}:
	\[
	\gamma = \alpha\circ(\beta_1,\dots,\beta_k) \in \mathcal{O}_G(n_1+\cdots+n_k)
	\]
	by
	\[
	\gamma(x_{1,1},\dots,x_{1,n_1};\dots;x_{k,1},\dots,x_{k,n_k})
	= \alpha\big(\beta_1(x_{1,1},\dots,x_{1,n_1}),\dots,\beta_k(x_{k,1},\dots,x_{k,n_k})\big).
	\]
	
	\begin{proposition}
		The substitution $\circ$ above defines a well-defined operation
		\[
		\mathcal{O}_G(k)\times \mathcal{O}_G(n_1)\times\cdots\times\mathcal{O}_G(n_k)
		\to \mathcal{O}_G(n_1+\cdots+n_k),
		\]
		compatible with the $G$-equivariance structure of twin operations.
	\end{proposition}
	
	\begin{proof}
		We must show that if $\alpha$ and each $\beta_i$ admit twin families $(\alpha_g)$ and $(\beta_{i,g})$,
		then $\gamma$ admits a twin family $(\gamma_g)$ given by the same transport rule, i.e.
		\[
		\gamma_g(\mathbf{x}) = g^{-1}\cdot \gamma(g\cdot\mathbf{x}),
		\]
		with $\mathbf{x}$ denoting the whole tuple of variables.
		Concretely, transport distributes over substitution. Using
		$\alpha_g(y_1,\dots,y_k)=g^{-1}\cdot\alpha(g\cdot y_1,\dots,g\cdot y_k)$ and
		$\beta_{i,g}(\mathbf{x}_i)=g^{-1}\cdot\beta_i(g\cdot\mathbf{x}_i)$, we compute
		\[
			\alpha_g\big(\beta_{1,g}(\mathbf{x}_1),\dots,\beta_{k,g}(\mathbf{x}_k)\big)
			= g^{-1}\cdot\alpha\big(g\cdot\beta_{1,g}(\mathbf{x}_1),\dots\big)
			= g^{-1}\cdot\alpha\big(\beta_1(g\cdot\mathbf{x}_1),\dots,\beta_k(g\cdot\mathbf{x}_k)\big)
			= g^{-1}\cdot\gamma(g\cdot\mathbf{x}) = \gamma_g(\mathbf{x}),
		\]
		using $g\cdot\beta_{i,g}(\mathbf{x}_i)=\beta_i(g\cdot\mathbf{x}_i)$ in the middle step. Hence
		$\gamma_g=\alpha_g\circ(\beta_{1,g},\dots,\beta_{k,g})$: the transport $\alpha\mapsto\alpha_g$ is an operad
		endomorphism of $\End_E$, so a substitution of twin operations is again a twin operation.
	\end{proof}
	
	\subsection{Symmetries and the symmetric operad structure}
	
	If $G$ acts on $E$ and the symmetric group $S_n$ acts on $E^n$ by permuting factors, we can also
	consider the interaction of twins with these symmetric group actions.
	
	\begin{definition}
		For $\sigma\in S_n$ and $\alpha\in\mathcal{O}_G(n)$, define
		\[
		(\alpha\cdot\sigma)(x_1,\dots,x_n)
		:= \alpha(x_{\sigma^{-1}(1)},\dots,x_{\sigma^{-1}(n)}).
		\]
	\end{definition}
	
	\begin{proposition}
		The collection $(\mathcal{O}_G(n))_{n\ge 1}$ with the substitution maps and the $S_n$-actions defined above
		forms a (symmetric) operad in the sense of~\cite{May1972,LodayVallette2012}.
	\end{proposition}
	
	\begin{proof}
		The axioms of a symmetric operad (associativity of substitution, unit, equivariance with respect
		to symmetric groups) are checked as in the standard example of the endomorphism operad of $E$,
		with the further observation that the twin structure is preserved because it is defined by transport
		along the group action, which commutes with the $S_n$-action.
	\end{proof}
	
	\begin{definition}[Twin operad]
		We call $\mathcal{O}_G$ the \emph{twin operad} associated with the $G$-object $E$.
	\end{definition}
	
	%===========================================================
	\section{Algebras over Twin Operads and Classical Structures}
	\label{sec:algebras}

	\begin{example}[Small cases of the twin operad]\label{ex:small-operad}
		When $G=\{e\}$ the action is trivial and $\mathcal{O}_{\{e\}}=\End_E$ is the ordinary endomorphism operad; a
		chosen associative $\alpha$ then determines an operad morphism $\mathsf{Assoc}\to\End_E$, as used below. For
		finite $E$ with $\abs{E}=m$ each arity is finite, $\abs{\mathcal{O}_G(n)}\le m^{m^{\,n}}$; for instance with
		$E=\Z/2$ there are $m^{m^2}=2^4=16$ binary operations, among which the unital associative commutative ones
		realise $E$ as a monoid. The twin operad refines this finite set by recording how $G$ permutes the operations
		through transport.
	\end{example}
	
	We now explain how familiar objects appear as algebras over $\mathcal{O}_G$.
	
	\subsection{Twin operad algebras}
	
	\begin{definition}
		Let $\mathcal{O}$ be an operad in a symmetric monoidal category $\mathcal{C}$.
		An \emph{$\mathcal{O}$-algebra} is an object $X\in\mathcal{C}$ together with structure maps
		\[
		\mathcal{O}(n)\otimes X^{\otimes n} \to X
		\]
		satisfying the usual associativity, unit, and equivariance axioms.
	\end{definition}
	
	\begin{definition}[Algebras over the twin operad]
		An \emph{$\mathcal{O}_G$-algebra} is an object $X$ of the ambient category equipped with a morphism of operads
		\[
		\mathcal{O}_G \to \End_X,
		\]
		where $\End_X(n)=\Hom(X^{\otimes n},X)$ is the endomorphism operad of $X$.
	\end{definition}
	
	\begin{remark}
		Concretely, an $\mathcal{O}_G$-algebra is a $G$-object $X$ endowed with twin operations
		acting on $X$ in a way compatible with operadic composition.
	\end{remark}
	
	\subsection{Associative $G$-equivariant algebras}
	
	Let $G$ act by automorphisms on an associative algebra $A$.
	
	\begin{proposition}
		If the action of $G$ preserves the associative multiplication of $A$, then $A$ is naturally an
		algebra over a quotient of $\mathcal{O}_G$ corresponding to the associative operad $\mathsf{Assoc}$.
	\end{proposition}
	
	\begin{proof}
		The associative operad $\mathsf{Assoc}$ is generated by a single binary operation $\mu$ modulo
		associativity, with $\mathsf{Assoc}(n)\cong k[S_n]$ the $n!$ orderings of the inputs. That $G$ preserves the
		multiplication of $A$ means $\mu$ lies in the $\LinearSubgroup$-fixed part of $\mathcal{O}_G(2)$, i.e.\ it is
		a genuine $G$-equivariant (twin) operation. A $G$-equivariant associative algebra structure 
		on $A$ can be encoded as a morphism of operads
		\[
		\mathsf{Assoc}\to\mathcal{O}_G\to\End_A,
		\]
		where the first arrow picks out the twin lift of the associative operation, and the second arrow
		evaluates it on $A$. The $G$-equivariance ensures compatibility with the twin structure.
	\end{proof}
	
	\subsection{Other structures}
	
	Similarly, by choosing appropriate quotients or sub-operads of $\mathcal{O}_G$, one can encode:
	
	\begin{itemize}
		\item commutative $G$-equivariant algebras via a commutative quotient of $\mathcal{O}_G$;
		\item Lie-type twin structures via an operad obtained from $\mathcal{O}_G$ by imposing antisymmetry and Jacobi-type identities;
		\item module structures and crossed products via colored or multi-sorted versions of $\mathcal{O}_G$.
	\end{itemize}
	
	We do not pursue these constructions further here, as they require a case-by-case analysis 
	depending on the type of algebraic identities involved.
	
	%===========================================================
	\section{Functoriality and 2-Functors}
	\label{sec:functoriality}
	
	The 2-category $\Twin$ comes with several natural 2-functors to more classical categories.
	
	\subsection{The forgetful 2-functor to $G$-objects}
	
	\begin{definition}
		Define a 2-functor
		\[
		U:\Twin \to \mathsf{GObj},
		\]
		where $\mathsf{GObj}$ is the 2-category of pairs $(E,G)$ with group actions, 1-morphisms
		given by equivariant maps, and 2-morphisms given by natural transformations (when enriched).
		On objects, $U(E,\alpha,G)=(E,G)$; on morphisms, $U(f,\phi)=(f,\phi)$; and on 2-morphisms,
		$U(\eta)=\eta$.
	\end{definition}
	
	\begin{proposition}
		$U$ is a strict 2-functor that is faithful on 1-morphisms and 2-morphisms.
	\end{proposition}
	
	\begin{proof}
		The functor $U$ simply forgets the twin operation $\alpha$; all coherence data for a strict
		2-functor are inherited from the structure of $\Twin$ and $\mathsf{GObj}$.
		Faithfulness follows from the fact that different 1- or 2-morphisms differ already at the level
		of the underlying $G$-objects.
	\end{proof}
	
	\subsection{The twin-orbit 2-functor}
	
	\begin{definition}
		Define a 2-functor
		\[
		\mathcal{O} : \Twin \to \mathsf{HomogeneousSpaces},
		\]
		sending $\TwinSpace$ to the homogeneous space $\LinearSubgroup\backslash G$, and a 1-morphism
		$F=(f,\phi)$ to the induced map of homogeneous spaces
		\[
		\LinearSubgroup\backslash G \to \LinearSubgroup'\backslash G',\qquad \LinearSubgroup g \mapsto \LinearSubgroup'\phi(g).
		\]
	\end{definition}
	
	\begin{proposition}
		$\mathcal{O}$ is a well-defined $2$-functor on the sub-$2$-category $\Twin^{\mathrm{epi}}$ whose
		$1$-morphisms have surjective underlying map $f$.
	\end{proposition}
	
	\begin{proof}
		Well-definedness on objects is \cref{prop:orbit}. For $1$-morphisms one must check $\phi(\LinearSubgroup)\subseteq\LinearSubgroup'$. Let
		$g\in\LinearSubgroup$, so $\alpha(g\cdot x,g\cdot y)=g\cdot\alpha(x,y)$. For $x'=f(x),\,y'=f(y)$ in the image of $f$,
		\[
			\alpha'(\phi(g)\cdot x',\phi(g)\cdot y')=\alpha'\big(f(g\cdot x),f(g\cdot y)\big)=f\big(\alpha(g\cdot x,g\cdot y)\big)
			=f\big(g\cdot\alpha(x,y)\big)=\phi(g)\cdot\alpha'(x',y').
		\]
		Since $f$ is surjective this holds for all $x',y'\in E'$, hence $\phi(g)\in\LinearSubgroup'$ and
		$\phi(\LinearSubgroup)\subseteq\LinearSubgroup'$; the induced coset map
		$\LinearSubgroup g\mapsto\LinearSubgroup'\phi(g)$ is therefore well-defined. (On a non-surjective image
		the relation $\phi(g)\in\LinearSubgroup'$ may fail, which is why the hypothesis is needed.) On 2-morphisms, the effect is trivial at the level
		of homogeneous spaces, so we obtain a strict 2-functor.
	\end{proof}
	
	\begin{remark}
		The 2-functor $\mathcal{O}$ encodes the geometric content of the twin structure: it ignores
		the detailed algebraic operation and remembers only the orbit $\LinearSubgroup\backslash G$.
		This reflects the central role of this homogeneous space in the deformation theory of twin spaces.
	\end{remark}
	
	%===========================================================
	\section{Interaction with Deformation Cohomology}
	\label{sec:interaction}
	
	We briefly explain how the deformation cohomology of Article~A can be viewed from this 2-categorical
	and operadic perspective.
	
	\subsection{From cochains to 2-endomorphisms}
	
	In Article~A, the deformation complex 
	\[
	\Cochain^\bullet(G,\End(E))
	\]
	was constructed as a Hochschild-type complex controlling deformations of $\TwinSpace$.
	
	\begin{remark}
		At a conceptual level, $\Cochain^n(G,\End(E))$ can be thought of as $n$-fold higher endomorphisms
		of the identity 1-morphism on $\TwinSpace$ in the 2-category $\Twin$, enriched appropriately.
		The differential $\diff$ encodes the failure of these higher endomorphisms to be compatible under
		composition, and the cohomology $\HH^\bullet$ measures higher deformations and obstructions to strictness.
		This viewpoint parallels the interpretation of Hochschild cohomology as self-extensions in the derived
		category of bimodules~\cite{Weibel1994}.
	\end{remark}
	
	\subsection{Operadic interpretation of $\HH^1$ and $\HH^2$}
	
	The twin operad $\mathcal{O}_G$ controls the allowed types of operations on $E$. A deformation of the
	twin structure can be seen as a deformation of the operad morphism
	\[
	\mathcal{O}_G \to \End_E.
	\]
	
	\begin{proposition}[Heuristic]
		Infinitesimal deformations parametrized by $\HH^1(G,\End(E))$ correspond to infinitesimal perturbations
		of the operad morphism $\mathcal{O}_G\to\End_E$, and obstructions in $\HH^2(G,\End(E))$ correspond to failures
		of extending such perturbations to higher order.
	\end{proposition}
	
	\begin{proof}[Idea of proof]
		Operadic deformation theory (as in the case of associative algebras) can be formulated in terms of Maurer--Cartan 
		elements in suitable differential graded Lie algebras built from the endomorphism operad.
		In the twin setting, the group action intertwines with the operadic structure, and the deformation 
		complex acquires the form $\Cochain^\bullet(G,\End(E))$. The compatibility between operadic and group-theoretic 
		structures yields the claimed identification, though a full proof would require setting up the operadic deformation
		complex explicitly, which we leave for future work.
	\end{proof}
	
	%===========================================================
	\section{Examples}
	\label{sec:examples}
	
	We conclude with a few examples illustrating the constructions above.
	
	\begin{example}[Matrix twin spaces]
		Let $E=M_n(k)$ with the conjugation action of $G=GL_n(k)$, and let $\alpha(X,Y)=XY$ be matrix multiplication.
		Then $(E,\alpha,G)$ is a twin space (as in Article~A), and 1-morphisms in $\Twin$ from this object to itself
		include inner automorphisms of $M_n(k)$ and certain outer automorphisms induced by group automorphisms of $GL_n(k)$.
		
		The twin operad $\mathcal{O}_G$ on $E$ includes all $G$-equivariant $n$-ary multilinear operations on $M_n(k)$.
		Algebras over suitable quotients of $\mathcal{O}_G$ recover associative $G$-equivariant algebra structures.
	\end{example}
	
	\begin{example}[Compact Lie group actions on manifolds]
		Let $G$ be a compact connected Lie group acting smoothly on a manifold $E$, and let $\alpha$ be a smooth
		$G$-equivariant operation (e.g.\ a pointwise product of functions, or a convolution-type operation).
		Then $(E,\alpha,G)$ is a twin space in the smooth category, yielding an object of $\Twin$.
		The monoidal structure of \cref{sec:monoidal} allows us to construct twin spaces on products
		of manifolds and group actions.
	\end{example}
	
	\begin{example}[Discretized twin operads]
		If $E$ is a finite $G$-set, then $\mathcal{O}_G(n)$ is a finite set of $G$-equivariant $n$-ary operations.
		The twin operad $\mathcal{O}_G$ is thus a finite operad in the category of sets.
		Algebras over $\mathcal{O}_G$ in the category of sets or vector spaces give discrete or linear models of
		twin structures, suitable for combinatorial or computational applications.
	\end{example}
	
	%===========================================================
	\section{Conclusion and Outlook}
	\label{sec:conclusion}
	
	In this article we have shown that twin spaces naturally form a 2-category $\Twin$ with a monoidal structure,
	and that operations on twin spaces can be organized into an operad $\mathcal{O}_G$ whose algebras encode
	a wide range of $G$-equivariant algebraic structures.
	
	From a conceptual perspective, the main outcome is that twin spaces live at the intersection of:
	\begin{itemize}
		\item group actions and homogeneous spaces ($\LinearSubgroup\backslash G$);
		\item higher category theory (via the 2-category $\Twin$ and its monoidal structure);
		\item operad theory (via the operads of twin operations and their algebras).
	\end{itemize}
	This viewpoint complements and enriches the deformation cohomology theory of Article~A, suggesting that
	twin spaces should be treated as higher-structured objects whose deformations are naturally approached
	via 2-categorical and operadic techniques.
	
	Possible directions for future work include:
	\begin{enumerate}[label=(\alph*)]
		\item \textbf{Enriched twin spaces.}
		Developing twin spaces in enriched categories (e.g.\ dg-categories, $\infty$-categories) so that both
		$\Twin$ and $\mathcal{O}_G$ become higher-categorical objects, suitable for derived deformation theory.
		\item \textbf{Homotopical twin operads.}
		Constructing model category structures on categories of twin operads and their algebras, allowing for 
		homotopy invariance and derived deformation theory in the sense of Lurie and others.
		\item \textbf{Applications to quantum groups and Hopf algebras.}
		Interpreting quantized structures as deformations of twin operads or as algebras in $\Twin$ equipped with
		additional duality and antipode data.
		\item \textbf{Links with higher gauge theory.}
		Investigating whether twin spaces and their 2-categorical structure can model aspects of higher gauge theories,
		where fields and symmetries are naturally described by higher groupoids and 2-groups.
	\end{enumerate}
	
	These directions suggest that twin spaces are not only a tool for organizing deformations of operations under 
	group actions, but also a promising building block for a broader higher-structural framework at the interface 
	of algebra, geometry, and mathematical physics.
	
	%===========================================================
	\begin{thebibliography}{99}
		
		\bibitem{Gerstenhaber1964}
		M.~Gerstenhaber,
		\newblock On the deformation of rings and algebras,
		\newblock {\em Ann. of Math. (2)} \textbf{79} (1964), 59--103.
		
		\bibitem{MacLane1998}
		S.~Mac Lane,
		\newblock {\em Categories for the Working Mathematician},
		\newblock 2nd ed., Graduate Texts in Mathematics, 5, Springer, 1998.
		
		\bibitem{Leinster2004}
		T.~Leinster,
		\newblock {\em Higher Operads, Higher Categories},
		\newblock London Math. Soc. Lecture Note Ser., 298, Cambridge Univ. Press, 2004.
		
		\bibitem{Kelly1982}
		G.~M.~Kelly,
		\newblock {\em Basic Concepts of Enriched Category Theory},
		\newblock London Math. Soc. Lecture Note Ser., 64, Cambridge Univ. Press, 1982.
		
		\bibitem{May1972}
		J.~P.~May,
		\newblock {\em The Geometry of Iterated Loop Spaces},
		\newblock Lecture Notes in Mathematics, 271, Springer, 1972.
		
		\bibitem{LodayVallette2012}
		J.-L.~Loday and B.~Vallette,
		\newblock {\em Algebraic Operads},
		\newblock Grundlehren der Mathematischen Wissenschaften, 346, Springer, 2012.
		
		\bibitem{Weibel1994}
		C.~A.~Weibel,
		\newblock {\em An Introduction to Homological Algebra},
		\newblock Cambridge Studies in Advanced Mathematics, 38, Cambridge Univ. Press, 1994.
		
	\end{thebibliography}
	
\end{document}
